% =====================================================================
% Paper completo -- MICCAI / Springer LNCS (español). Requiere llncs.cls + babel.
% Diagramas: figures/figNN_es.pdf. Marcadores: [[ ... ]]
% =====================================================================
\documentclass[runningheads]{llncs}
\usepackage[utf8]{inputenc}
\usepackage[T1]{fontenc}
\usepackage[spanish,es-noshorthands]{babel}
\usepackage{graphicx}
\usepackage{booktabs}
\usepackage{amsmath,amssymb}
\usepackage{multirow}
\usepackage{xcolor}
\usepackage[hidelinks]{hyperref}
\renewcommand{\abstractname}{Resumen}
\renewcommand{\keywordname}{\bfseries Palabras clave:}
\newcommand{\todo}[1]{\textcolor{red}{[[#1]]}}
\newcommand{\dfig}[3]{\begin{figure}[t]\centering
\includegraphics[width=\textwidth,height=0.34\textheight,keepaspectratio]{#1}
\caption{#2}\label{#3}\end{figure}}

\title{Redes Convolucionales Cuánticas frente a Clásicas para la Segmentación\\
de Subregiones del Tronco Encefálico en la Ataxia Espinocerebelosa\\
Tipo~2: Un Estudio Comparativo Integral con Parámetros Equiparados}
\titlerunning{Redes Cuánticas vs Clásicas para Segmentación del Tronco Encefálico (SCA2)}
\author{\todo{Autor Uno}\inst{1} \and \todo{Autor Dos}\inst{1} \and \todo{Autor Senior}\inst{1}}
\authorrunning{\todo{Autor Uno} et al.}
\institute{Centro de Neurociencias de Cuba (CNEURO), La Habana, Cuba\\
\email{\todo{correspondencia@email}}}

\begin{document}
\maketitle

\begin{abstract}
La segmentación precisa de las subregiones del tronco encefálico (bulbo
raquídeo, protuberancia, mesencéfalo) sustenta los biomarcadores de imagen de la
neurodegeneración, en particular la atrofia progresiva de la protuberancia y el
bulbo característica de la ataxia espinocerebelosa tipo~2 (SCA2), un trastorno
cuya mayor prevalencia mundial se registra en Holguín, Cuba. Presentamos una
comparación integral y controlada de nueve arquitecturas de segmentación sobre
una cohorte relevante para la SCA2 adquirida en el Centro de Neurociencias de
Cuba, enfrentando seis U-Nets clásicas a tres redes híbridas cuántico--clásicas,
cada una en 2D y 3D con un presupuesto equiparado de ${\sim}17$M parámetros.
Todos los modelos alcanzan precisión elevada (Dice medio $0{,}92$--$0{,}94$). Una
ablación de capacidad sitúa el techo de los modelos cuánticos en el cuello de
botella cuántico (alcanzan la precisión clásica con ${\sim}\tfrac13$ de los
parámetros, a un costo computacional simulado un orden de magnitud mayor). La
metodología de evaluación domina las conclusiones en 2D: bajo la métrica
convencional todos los modelos 2D son indistinguibles (Friedman $p{=}0{,}637$),
pero bajo evaluación honesta sobre el volumen completo se separan con claridad
($p{<}0{,}0001$), siendo los modelos híbridos cuánticos los más robustos frente a
falsos positivos en cortes vacíos (agrupado $p{=}0{,}031$); en 3D los modelos
clásicos conservan una pequeña ventaja ($p{=}0{,}008$). Aportamos diagramas por
arquitectura y delineamos líneas futuras que incluyen hardware cuántico real,
validación multicéntrica y el uso de mapas de atrofia poblacional, mapas
autoorganizados y mapas de incertidumbre. Todos los procedimientos siguieron la
Declaración de Helsinki.
\keywords{Segmentación del tronco encefálico \and Ataxia espinocerebelosa \and
Aprendizaje automático cuántico \and U-Net \and Metodología de evaluación.}
\end{abstract}

\section{Introducción}
La ataxia espinocerebelosa tipo~2 (SCA2) es un trastorno neurodegenerativo
autosómico dominante causado por una expansión de repeticiones CAG en
\emph{ATXN2}~\cite{pulst}. Su neuropatología se centra en el sistema
olivopontocerebeloso; la atrofia cuantitativa de la protuberancia y el bulbo es
un biomarcador de imagen establecido y sensible a la progresión~\cite{velazquez}.
La SCA2 reviste singular importancia en Cuba, donde Holguín presenta la mayor
prevalencia del mundo y donde el Centro de Neurociencias de Cuba (CNEURO)
mantiene un programa de décadas~\cite{velazquez}. Derivar estos biomarcadores a
gran escala exige una delineación automática y fiable; el trazado manual es lento
y dependiente del observador.

Las U-Nets convolucionales~\cite{unet,unet3d} y numerosas variantes
\cite{cerebnet,acapulco,swinunetr,inception,cbam} dominan la segmentación médica.
Las redes híbridas cuántico--clásicas se han propuesto como aproximadores
eficientes en parámetros~\cite{qcnn}, pero su utilidad para la segmentación
volumétrica apenas se ha probado y rara vez se compara a igualdad de capacidad
bajo un protocolo clínicamente honesto. Planteamos: (i)~a igualdad de capacidad,
¿difieren los modelos clásicos y cuánticos? (ii)~¿dónde está su techo? (iii)~¿las
conclusiones dependen del protocolo de evaluación? Nuestras contribuciones son
una comparación reproducible de nueve modelos en 2D y 3D con diagramas por
arquitectura, una ablación de capacidad, la evidencia de que la evaluación
honesta sobre el volumen completo revierte la aparente equivalencia 2D, y una
hoja de ruta que incluye el uso de mapas en estudios futuros.

\section{Trabajo Relacionado}
La U-Net~\cite{unet} y su forma 3D~\cite{unet3d} establecieron el paradigma
codificador--decodificador con conexiones de salto; ACAPULCO~\cite{acapulco} y
CerebNet~\cite{cerebnet} lo adaptaron al cerebelo/tronco. Los codificadores
transformer~\cite{swin,swinunetr} añadieron atención de largo alcance;
inception~\cite{inception} y la atención por compuertas~\cite{cbam} mejoraron el
modelado multiescala y de saltos. Las CNN cuánticas~\cite{qcnn} y los circuitos
diferenciables~\cite{pennylane,qiskit} se han probado sobre todo en tareas de
baja dimensión; las comparaciones justas y equiparadas en datos clínicos
volumétricos son escasas. Por último, los flujos 2D que filtran a cortes con
estructura pueden inflar la precisión aparente, algo que cuantificamos.

\section{Datos: Imagen Cerebral y RM}
\subsection{Cohorte y ética}
Se adquirieron RM cerebrales potenciadas en T1 de \todo{$N$} sujetos
(\todo{edad, sexo, composición clínica}) en CNEURO. \textbf{Todos los
procedimientos se ajustaron a la Declaración de Helsinki}~\cite{helsinki}; el
estudio fue aprobado por \todo{comité de ética, n.\ de aprobación} con
consentimiento informado por escrito.

\subsection{Adquisición y preprocesamiento}
Los volúmenes se procesaron con corrección de campo de sesgo N4~\cite{n4},
registro afín al espacio MNI, recorte a un FOV de $80\times80\times96$ centrado
en el tronco y normalización z por volumen (sin estadísticas entre sujetos;
Fig.~\ref{fig:prep}). Se etiquetan cuatro clases: fondo, bulbo raquídeo,
protuberancia y mesencéfalo.
\dfig{figures/fig01_es.pdf}{Flujo de preprocesamiento de RM.}{fig:prep}

\subsection{Particiones de datos}
Los 39 sujetos se dividen a nivel de sujeto (semilla~42): 27 entrenamiento / 6
validación / 6 prueba, con un conjunto de prueba idéntico para todos los modelos
(Fig.~\ref{fig:split}).
\dfig{figures/fig02_es.pdf}{Partición de datos a nivel de sujeto.}{fig:split}

\section{Métodos: Arquitecturas}
Todos los modelos son agnósticos a la dimensión (un indicador selecciona
operadores 2D/3D) y comparten una plantilla codificador--decodificador con
saltos (Fig.~\ref{fig:unet}).
\dfig{figures/fig03_es.pdf}{Plantilla U-Net compartida.}{fig:unet}

\subsection{ClassicalCNN}
U-Net de doble convolución (\texttt{Conv-BatchNorm-ReLU}$\times2$ por etapa),
canales $[64,128,256,352]$ (2D)/$[38,76,156,204]$ (3D), cuatro submuestreos,
saltos por concatenación; ${\sim}16{,}95$M (Fig.~\ref{fig:ccnn}).
\dfig{figures/fig04_es.pdf}{ClassicalCNN.}{fig:ccnn}

\subsection{CerebNet}
Bloques densos competitivos (\texttt{Conv-BatchNorm-PReLU} con atajo de maxout
competitivo), MaxPool/MaxUnpool por índices, $160$ canales uniformes;
${\sim}16{,}65$M (2D)/$17{,}98$M (3D) (Fig.~\ref{fig:cereb}).
\dfig{figures/fig05_es.pdf}{CerebNet.}{fig:cereb}

\subsection{MipaimUNet}
Cinco niveles de codificación con bloques inception multi-rama~\cite{inception}
con InstanceNorm, saltos refinados por CBAM~\cite{cbam} y cuello de botella
inception explícito; ${\sim}17{,}27$M (2D)/$17{,}80$M (3D) (Fig.~\ref{fig:mip}).
\dfig{figures/fig06_es.pdf}{MipaimUNet.}{fig:mip}

\subsection{SwinUNETR}
Codificador transformer Swin~\cite{swin,swinunetr} (atención por ventanas
desplazadas, ventana~7) y decodificador convolucional con saltos por
concatenación; ${\sim}16{,}79$M (2D)/$17{,}06$M (3D) (Fig.~\ref{fig:swin}).
\dfig{figures/fig07_es.pdf}{SwinUNETR.}{fig:swin}

\subsection{ACAPULCO}
U-Net residual con preactivación (\texttt{InstanceNorm-ReLU-Conv} con suma
residual), dos bloques por etapa, canales $[44,88,176,240]$ (2D);
${\sim}17{,}38$M (2D)/$17{,}08$M (3D) (Fig.~\ref{fig:aca}).
\dfig{figures/fig08_es.pdf}{ACAPULCO.}{fig:aca}

\subsection{MARIN}
Diseño de fusión: ordenación con preactivación, saltos con compuerta CBAM, cuello
de botella con maxout competitivo y ramas dobles dilatadas (dilatación~1~$\|$~2)
por bloque, con PReLU e InstanceNorm; ${\sim}17{,}42$M (2D)/$17{,}10$M (3D)
(Fig.~\ref{fig:marin}).
\dfig{figures/fig09_es.pdf}{MARIN.}{fig:marin}

\subsection{Redes híbridas cuántico--clásicas}
Los tres modelos ``cuánticos'' comparten un diseño. Una U-Net convolucional
clásica (codificador de 3 etapas $[80,160,320]$) reduce su cuello de botella $x$
por promediado global a un vector de canales (se descarta la estructura
espacial), lo pasa por un circuito cuántico variacional (Fig.~\ref{fig:circ}), lo
reproyecta, lo difunde y lo \emph{suma} a $x$ ($x \leftarrow x + q_{\text{feat}}$).
La vía cuántica es una rama lateral global de recalibración de canales
(Fig.~\ref{fig:hyb}), no un cuello de botella por el que fluya la representación;
las características convolucionales pasan inalteradas. PennyLane-QCNN/Hybrid-QCNN
usan un circuito PennyLane (codificación de ángulos + StronglyEntanglingLayers, 8
qubits); Qiskit-QCNN usa un EstimatorQNN de Qiskit (ZZFeatureMap + RealAmplitudes,
4 qubits). \textbf{Ninguno es puramente cuántico; todos son híbridos.} Los
circuitos se simulan de forma clásica.
\dfig{figures/fig10_es.pdf}{Arquitectura híbrida cuántico--clásica (rama
cuántica aditiva).}{fig:hyb}
\dfig{figures/fig11_es.pdf}{Circuito cuántico variacional del cuello de
botella.}{fig:circ}

\subsection{Equiparación de parámetros}
Los modelos clásicos mantienen ${\sim}17$M parámetros en 2D y 3D. El codificador
fijo de 3 etapas de los modelos cuánticos da ${\sim}6{,}1$M en 2D pero
${\sim}17{,}6$M en 3D (los núcleos 3D cuestan ${\sim}3\times$), por lo que los
modelos cuánticos 2D nativos no están equiparados. Añadimos una ejecución 2D con
codificador de 4 etapas $[80,160,320,480]\,({=}16{,}9$M, variantes ``17M''). El
circuito en sí aporta parámetros despreciables.

\section{Métodos: Entrenamiento, Evaluación, Hardware}
Los modelos se entrenaron con Adam~\cite{adam} (tasa~$10^{-3}$, coseno,
calentamiento de 5 ep), una pérdida combinada de entropía cruzada~+~Dice
suave~\cite{vnet} con ponderación de clases por frecuencia inversa, precisión
mixta, recorte de gradiente, parada temprana (paciencia~15, mínimo 250 épocas),
hasta 400 épocas, lote 16 (2D)/2 (3D), en una RTX~3090.

\paragraph{Métricas y protocolo.} Reportamos Dice por clase y medio de primer
plano. La métrica de entrenamiento es un micro-promedio; para la estadística
reportamos la media del Dice por sujeto~$\pm$~d.e.\ ($n{=}6$). Como el filtro de
20 vóxeles se aplica al conjunto de prueba 2D, reportamos además una evaluación
honesta sobre el volumen completo en 2D sobre los 96 cortes
(Fig.~\ref{fig:eval}). Estadística: IC del 95\% por remuestreo de diferencias
pareadas y $d$ de Cohen como primarias; Friedman y Wilcoxon corregido por Holm
como exploratorias.
\dfig{figures/fig12_es.pdf}{Protocolo de evaluación.}{fig:eval}

\paragraph{Software.} Python~3.12, PyTorch~2.6 (CUDA~12.4), PennyLane~0.44,
Qiskit~2.4; $2{\times}$RTX~3090 (24\,GB). El flujo completo se muestra en la
Fig.~\ref{fig:flow}.
\dfig{figures/fig13_es.pdf}{Flujo experimental completo.}{fig:flow}

\section{Resultados}
Todos los valores de Dice son media por sujeto~$\pm$~d.e.\ ($n{=}6$).

\subsection{Precisión 2D y 3D a igualdad de capacidad}
Sobre cortes 2D con estructura los nueve modelos fueron indistinguibles (Friedman
$\chi^2{=}6{,}09$, $p{=}0{,}637$); en 3D difirieron ($\chi^2{=}20{,}58$,
$p{=}0{,}008$). El mesencéfalo es la subregión más difícil y el bulbo la más
fácil (Tablas~\ref{tab:2d},~\ref{tab:3d}).

\begin{table}[t]\centering
\caption{2D, cortes con estructura. Q/H/C = cuántico/híbrido/clásico.}
\label{tab:2d}\small\setlength{\tabcolsep}{5pt}
\begin{tabular}{llccccc}
\toprule
Modelo & T & Par. & Dice medio & bulbo & protub. & mesenc.\\
\midrule
Qiskit-QCNN & Q & 16,9M & $0{,}932{\pm}0{,}007$ & 0,947 & 0,935 & 0,915\\
MARIN & C & 17,4M & $0{,}932{\pm}0{,}015$ & 0,948 & 0,939 & 0,909\\
Hybrid-QCNN & H & 16,9M & $0{,}932{\pm}0{,}011$ & 0,946 & 0,938 & 0,910\\
MipaimUNet & C & 17,3M & $0{,}931{\pm}0{,}015$ & 0,948 & 0,936 & 0,910\\
PennyLane-QCNN & Q & 16,9M & $0{,}930{\pm}0{,}013$ & 0,947 & 0,936 & 0,906\\
SwinUNETR & C & 16,8M & $0{,}929{\pm}0{,}014$ & 0,943 & 0,935 & 0,910\\
ClassicalCNN & C & 17,0M & $0{,}929{\pm}0{,}013$ & 0,944 & 0,934 & 0,908\\
ACAPULCO & C & 17,4M & $0{,}928{\pm}0{,}017$ & 0,936 & 0,933 & 0,915\\
CerebNet & C & 16,7M & $0{,}927{\pm}0{,}016$ & 0,937 & 0,929 & 0,915\\
\bottomrule
\end{tabular}
\end{table}

\begin{table}[t]\centering
\caption{3D, volúmenes completos.}
\label{tab:3d}\small\setlength{\tabcolsep}{5pt}
\begin{tabular}{llccccc}
\toprule
Modelo & T & Par. & Dice medio & bulbo & protub. & mesenc.\\
\midrule
MipaimUNet & C & 17,8M & $0{,}940{\pm}0{,}013$ & 0,943 & 0,942 & 0,933\\
ACAPULCO & C & 17,1M & $0{,}937{\pm}0{,}017$ & 0,938 & 0,942 & 0,930\\
PennyLane-QCNN & Q & 17,6M & $0{,}935{\pm}0{,}016$ & 0,943 & 0,942 & 0,921\\
MARIN & C & 17,1M & $0{,}934{\pm}0{,}019$ & 0,948 & 0,938 & 0,916\\
SwinUNETR & C & 17,1M & $0{,}932{\pm}0{,}020$ & 0,944 & 0,938 & 0,914\\
Hybrid-QCNN & H & 17,6M & $0{,}931{\pm}0{,}022$ & 0,935 & 0,935 & 0,923\\
CerebNet & C & 18,0M & $0{,}930{\pm}0{,}018$ & 0,939 & 0,937 & 0,915\\
Qiskit-QCNN & Q & 17,6M & $0{,}925{\pm}0{,}032$ & 0,919 & 0,934 & 0,921\\
ClassicalCNN & C & 17,0M & $0{,}920{\pm}0{,}037$ & 0,933 & 0,917 & 0,913\\
\bottomrule
\end{tabular}
\end{table}

\subsection{Ablación de capacidad}
Ampliar el codificador cuántico 2D $6{,}1$M${\to}16{,}9$M no dio mejora
(PennyLane $\Delta{=}{-}0{,}0023$, IC~95\% $[-0{,}0046,-0{,}0005]$; Qiskit
$-0{,}0007$; Hybrid $+0{,}0019$; ninguno significativo tras Holm). El cuello de
botella de 8 qubits, no la red troncal, limita el rendimiento.

\subsection{Evaluación honesta sobre el volumen completo}
Puntuar los 96 cortes de cada sujeto eleva Friedman de $p{=}0{,}637$ a
$\chi^2{=}39{,}5$, $p{<}0{,}0001$, ampliando la dispersión de ${\approx}0{,}006$ a
$0{,}059$ (Tabla~\ref{tab:c2}, Fig.~\ref{fig:c2}). El filtro infló el Dice 2D en
$0{,}018$--$0{,}087$ y reordenó el ranking (MARIN, co-líder, cae al último
puesto). Los modelos híbridos cuánticos son los más robustos; agrupado,
cuántico/híbrido vs clásico $0{,}894$ vs $0{,}878$, $\Delta{=}{+}0{,}016$,
Wilcoxon $p{=}0{,}031$.

\begin{table}[t]\centering
\caption{Dice 2D: solo primer plano vs volumen completo honesto ($n{=}6$).}
\label{tab:c2}\setlength{\tabcolsep}{6pt}
\begin{tabular}{llccc}
\toprule
Modelo & Tipo & Filtrado & Vol.\ completo & $\Delta$\\
\midrule
Qiskit-QCNN (17M) & Q & 0,932 & \textbf{0,905} & $-0,028$\\
CerebNet & C & 0,927 & 0,899 & $-0,028$\\
ClassicalCNN & C & 0,929 & 0,898 & $-0,031$\\
Hybrid-QCNN (17M) & H & 0,932 & 0,898 & $-0,034$\\
SwinUNETR & C & 0,929 & 0,883 & $-0,046$\\
PennyLane-QCNN (17M) & Q & 0,930 & 0,880 & $-0,050$\\
MipaimUNet & C & 0,931 & 0,875 & $-0,056$\\
ACAPULCO & C & 0,928 & 0,866 & $-0,062$\\
MARIN & C & 0,932 & \textbf{0,845} & $-0,087$\\
\midrule
PennyLane-QCNN (6M) & Q & 0,932 & 0,914 & $-0,018$\\
Qiskit-QCNN (6M) & Q & 0,933 & 0,912 & $-0,022$\\
Hybrid-QCNN (6M) & H & 0,930 & 0,899 & $-0,031$\\
\bottomrule
\end{tabular}
\end{table}

\begin{figure}[t]\centering
\includegraphics[width=0.82\textwidth]{figures/fig_fullvol_vs_fg_2d_es.png}
\caption{Dice 2D por sujeto: solo primer plano (gris) vs volumen completo (rojo =
cuántico/híbrido, azul = clásico), ordenado por la puntuación de volumen
completo.}
\label{fig:c2}
\end{figure}

\subsection{Costo computacional}
Las ejecuciones cuánticas 2D tardaron $8{,}9$--$13{,}6$\,h frente a
$0{,}3$--$1{,}0$\,h de los clásicos; en 3D la brecha se reduce (cuántico
$28$--$35$\,min). El costo refleja la simulación de vector de estado por muestra
y limitada por CPU (Fig.~\ref{fig:cost}).
\begin{figure}[t]\centering
\includegraphics[width=0.6\textwidth]{figures/fig_cost_2d_es.png}
\caption{Precisión 2D frente al tiempo de entrenamiento (escala log).}
\label{fig:cost}
\end{figure}

\section{Discusión}
A igualdad de capacidad, las U-Nets híbridas cuántico--clásicas igualan a las
clásicas en 2D y quedan ligeramente por detrás en 3D. Su ventaja es de eficiencia
en parámetros, no de precisión: un cuello de botella de 8 qubits sustituye a
varios millones de parámetros convolucionales sin pérdida de Dice, pero limita la
precisión alcanzable. El hallazgo más relevante es metodológico: la puntuación
honesta sobre el volumen completo separó modelos que el protocolo filtrado hacía
idénticos y reordenó el ranking, siendo los modelos híbridos cuánticos (y las
variantes estrechas de 6,1M) los más robustos frente a falsos positivos en cortes
vacíos ---posiblemente porque la rama cuántica global aditiva actúa como una
compuerta conservadora de canales. Los benchmarks de segmentación del tronco
encefálico deben, pues, reportarse sobre volúmenes completos. Clínicamente, todos
los modelos fueron lo bastante precisos (Dice de bulbo/protuberancia ${\sim}0{,}94$)
para sustentar mediciones de atrofia en SCA2; una U-Net 3D clásica y rápida es la
opción pragmática de despliegue.

\section{Trabajo Futuro}
\textbf{Metodología:} entrenamiento con múltiples semillas, cohortes
multicéntricas mayores (el $n{=}6$ actual limita la potencia), reentrenamiento con
volumen completo y con la pérdida de Dice ponderada corregida
(\S\ref{sec:lim}); métricas de borde (Hausdorff, Dice de superficie).
\textbf{Cuántico:} ejecución en hardware real con entrenamiento consciente del
ruido, codificaciones espacialmente conscientes más allá del promediado global, y
evaluación vectorizada del circuito. \textbf{Mapas en estudios futuros
(Fig.~\ref{fig:maps}):} (i)~\emph{atlas de atrofia/probabilísticos}
poblacionales ---estadística por vóxel en la cohorte para localizar la
degeneración pontina/medular y correlacionarla con escalas clínicas y la longitud
de repetición CAG, convirtiendo la segmentación en una tubería de biomarcadores;
(ii)~\emph{mapas autoorganizados}~\cite{som} y SOM cuánticos de las
representaciones aprendidas para el descubrimiento no supervisado de subtipos, una
extensión natural del tema cuántico; (iii)~\emph{mapas de saliencia e
incertidumbre} (MC-dropout/ensembles) para la interpretabilidad y la confianza
clínica.
\dfig{figures/fig14_es.pdf}{De las segmentaciones por sujeto a los mapas de
atrofia poblacional.}{fig:maps}

\section{Limitaciones}\label{sec:lim}
Conjunto de prueba pequeño ($n{=}6$; significancia con poca potencia, tamaños del
efecto primarios); una sola semilla; datos de un solo centro; circuitos cuánticos
simulados de forma clásica. Una auditoría del código halló que el término de Dice
ponderado por clases se multiplicaba por los pesos sin renormalizar, escalando mal
la pérdida (y produciendo los valores negativos observados durante el
entrenamiento); la corrección se aplica para ejecuciones futuras, por lo que la
comparación \emph{relativa} reportada es internamente consistente pero los valores
absolutos pueden variar tras reentrenar.

\section{Conclusión}
Sobre una cohorte de tronco encefálico relevante para la SCA2 del CNEURO, las
redes de segmentación cuánticas y clásicas con parámetros equiparados son
comparablemente precisas, ofreciendo los modelos cuánticos eficiencia en
parámetros a un costo computacional simulado sustancial. La evaluación honesta
sobre el volumen completo revela a los modelos híbridos cuánticos como los más
robustos frente a falsos positivos en cortes vacíos, mientras que en 3D los
clásicos conservan una pequeña ventaja. Los hallazgos respaldan la investigación
continua y controlada de la segmentación híbrida cuántico--clásica, los estándares
de evaluación sobre volúmenes completos y el uso de mapas de atrofia,
autoorganizados y de incertidumbre en estudios futuros.

\subsubsection*{Agradecimientos.} \todo{Financiación y agradecimientos de
recolección de datos.}

\begin{thebibliography}{99}
\bibitem{unet} Ronneberger, O., Fischer, P., Brox, T.: U-Net: convolutional networks for biomedical image segmentation. MICCAI (2015)
\bibitem{unet3d} \c{C}i\c{c}ek, \"O., et al.: 3D U-Net: learning dense volumetric segmentation from sparse annotation. MICCAI (2016)
\bibitem{vnet} Milletari, F., Navab, N., Ahmadi, S.A.: V-Net: fully convolutional neural networks for volumetric medical image segmentation. 3DV (2016)
\bibitem{cerebnet} Faber, J., et al.: CerebNet: a fast and reliable deep-learning pipeline for detailed cerebellum sub-segmentation. NeuroImage (2022)
\bibitem{acapulco} Han, S., Carass, A., et al.: Automatic cerebellum anatomical parcellation using U-Net with locally constrained optimization (ACAPULCO). NeuroImage (2020)
\bibitem{swinunetr} Hatamizadeh, A., et al.: Swin UNETR: Swin transformers for semantic segmentation of brain tumors in MRI images. MICCAI BrainLes (2022)
\bibitem{swin} Liu, Z., et al.: Swin Transformer: hierarchical vision transformer using shifted windows. ICCV (2021)
\bibitem{inception} Szegedy, C., et al.: Going deeper with convolutions. CVPR (2015)
\bibitem{cbam} Woo, S., Park, J., Lee, J.Y., Kweon, I.S.: CBAM: convolutional block attention module. ECCV (2018)
\bibitem{qcnn} Cong, I., Choi, S., Lukin, M.D.: Quantum convolutional neural networks. Nature Physics 15, 1273--1278 (2019)
\bibitem{pennylane} Bergholm, V., et al.: PennyLane: automatic differentiation of hybrid quantum-classical computations. arXiv:1811.04968 (2018)
\bibitem{qiskit} Qiskit contributors: Qiskit: an open-source framework for quantum computing (2023)
\bibitem{n4} Tustison, N.J., et al.: N4ITK: improved N3 bias correction. IEEE TMI 29(6), 1310--1320 (2010)
\bibitem{adam} Kingma, D.P., Ba, J.: Adam: a method for stochastic optimization. ICLR (2015)
\bibitem{pulst} Pulst, S.M., et al.: Moderate expansion of a normally biallelic trinucleotide repeat in spinocerebellar ataxia type 2. Nature Genetics 14, 269--276 (1996)
\bibitem{velazquez} Velazquez-Perez, L., Rodriguez-Labrada, R., Fernandez-Ruiz, J.: Spinocerebellar ataxia type 2: clinicogenetic aspects, mechanistic insights, and management. \todo{verificar revista/año}
\bibitem{som} Kohonen, T.: The self-organizing map. Proceedings of the IEEE 78(9), 1464--1480 (1990)
\bibitem{helsinki} Asociación Médica Mundial: Declaración de Helsinki: principios éticos para las investigaciones médicas en seres humanos. JAMA 310(20), 2191--2194 (2013)
\bibitem{nonparam} Friedman, M.: The use of ranks to avoid the assumption of normality. JASA 32, 675--701 (1937); Wilcoxon, F.: Individual comparisons by ranking methods. Biometrics Bulletin 1, 80--83 (1945)
\end{thebibliography}
\end{document}
